\subsection{Complexidade O(n²)}

A complexidade O(n²) é uma medida de desempenho de algoritmos que indica que o tempo de execução cresce quadraticamente com o tamanho da entrada, representado por $'n'$. Isso significa que, se o tamanho da entrada dobrar, o tempo de execução será aproximadamente quatro vezes maior, devido ao fator quadrático. Algoritmos com essa complexidade frequentemente envolvem dois loops aninhados, onde cada elemento da entrada é comparado ou processado em relação a todos os outros elementos.

Um exemplo clássico de algoritmo com complexidade O(n²) é o Bubble Sort, onde, para cada elemento de uma lista, são feitas comparações com os elementos subsequentes para realizar trocas e ordenar os dados. Em um caso médio ou pior, o algoritmo precisa percorrer a lista n vezes, e em cada iteração, realiza até n comparações ou operações, resultando em n * n = n² aoperações. Essa característica torna esses algoritmos ineficientes para grandes volumes de dados, pois o custo computacional aumenta rapidamente.

Embora algoritmos O(n²) sejam menos eficientes do que aqueles com complexidades lineares O(n) ou logarítmicas O(log n), eles ainda têm utilidade em cenários com entradas pequenas ou quando a simplicidade de implementação é mais importante que o desempenho. Além disso, podem servir como base para otimizações ou em situações específicas onde o pior caso raramente ocorre. No entanto, para aplicações de larga escala, é comum buscar alternativas mais eficientes, como algoritmos com complexidade O(n log n), típicos de métodos de ordenação mais avançados como Quick Sort ou Merge Sort.

Segue o exemplo representativo da complexidade O(n²):
\begin{center}
Fibonacci em Vetor

{Beecrowd 1176}

Timelimit: $1$	
\end{center}

Faça um programa que leia um valor e apresente o número de Fibonacci correspondente a este valor lido. Lembre que os 2 primeiros elementos da série de Fibonacci são 0 e 1 e cada próximo termo é a soma dos 2 anteriores a ele. Todos os valores de Fibonacci calculados neste problema devem caber em um inteiro de 64 bits sem sinal.

\vspace{0.3cm}
\textbf{Entrada}

A primeira linha da entrada contém um inteiro T, indicando o número de casos de teste. Cada caso de teste contém um único inteiro N $(0 \leq N \leq 60)$, correspondente ao N-esimo termo da série de Fibonacci.

\vspace{0.3cm}
\textbf{Saída}

Para cada caso de teste da entrada, imprima a mensagem "Fib(N) = X", onde X é o N-ésimo termo da série de Fibonacci.

%---------------------------------------------------------

\begin{table}[!ht]
    \centering
    \begin{tabular}{|p{7cm}|p{7cm}|}
        \hline
        \begin{center}
            \vspace{-0.3cm}
            \textbf{Exemplos de Entrada}
            \vspace{-0.3cm}
        \end{center} 
         &  
        \begin{center}
            \vspace{-0.3cm}
            \textbf{Exemplos de Saída}
            \vspace{-0.3cm}
        \end{center} \\ \hline
         \texttt{3} & \texttt{FIB(0) = 0} \\
         \texttt{0} & \texttt{FIB(4) = 3} \\
         \texttt{4} & \texttt{FIB(2) = 1} \\
         \texttt{2} & \texttt{}\\
         \hline
    \end{tabular}
    \caption{Exemplos de entradas e saídas}
    \label{tabela04}
\end{table}
%------------------------------------


\begin{center}
\begin{minipage}{0.5\textwidth}
    \lstinputlisting[language=C++]{fig/fibo.cpp}
\end{minipage}
\end{center}
$T(n)=C1*1+C2*1+C3*1+C4*(59)+C5*T+C6*T$\\
$T(n)=K+(C5+C6)*T$\\
$T(n)=a*T+b=O(T)$